% ===== C 卷参考答案与评分标准 =====
\anstou{C}

\datiA{名词解释（每小题 5 分，共 20 分）}

\thao{1.}\da 规律是事物运动发展中本质的、必然的、稳定的联系。\pingfen{3 分}规律具有客观性，不以人的意志为转移，但人可以认识和利用规律。\pingfen{2 分}

\thao{2.}\da 价值是揭示外部客观世界对于满足人的需要的意义关系的范畴，是客体属性对主体需要的满足关系。\pingfen{3 分}价值具有主体性、客观性和社会历史性，真理尺度与价值尺度在实践中的统一是成功实践的前提。\pingfen{2 分}

\thao{3.}\da 资本积累是把剩余价值再转化为资本，即剩余价值的资本化。\pingfen{3 分}其源泉是剩余价值，实质是用无偿占有的剩余价值扩大资本规模，从而加深两极分化。\pingfen{2 分}

\thao{4.}\da 共产主义是无产阶级解放和人类解放的社会形态，是物质财富极大丰富、人民精神境界极大提高、每个人自由而全面发展的社会。\pingfen{3 分}共产主义既是远大理想，也是现实运动，社会主义是它的必要准备和必经阶段。\pingfen{2 分}

\datiA{简答题（每小题 8 分，共 32 分）}

\thao{1.}\yaodian
\begin{enumerate}[label=(\arabic*),leftmargin=2.2em,itemsep=2pt,topsep=2pt]
\item 能动地认识世界：意识活动具有目的性、计划性、创造性和指导性，能反映事物的本质和规律。\pingfen{3 分}
\item 能动地改造世界：意识通过指导实践把观念的东西变为现实，对人体生理活动也有调控作用。\pingfen{3 分}
\item 正确发挥的条件：尊重客观规律是前提；必须通过实践这一根本途径；还要依赖一定的物质条件和物质手段。\pingfen{2 分}
\end{enumerate}

\thao{2.}\yaodian
\begin{enumerate}[label=(\arabic*),leftmargin=2.2em,itemsep=2pt,topsep=2pt]
\item 辩证否定是事物的自我否定，是事物内部矛盾运动的结果。\pingfen{2 分}
\item 辩证否定的实质是“扬弃”，即既克服又保留，克服旧事物中过时的东西，保留其中积极的合理的因素。\pingfen{3 分}
\item 辩证否定是发展环节和联系环节的统一。\pingfen{1 分}
\item 与形而上学否定的区别：形而上学否定是外力否定、全盘抛弃，割断了新旧事物之间的联系。\pingfen{2 分}
\end{enumerate}

\thao{3.}\yaodian
\begin{enumerate}[label=(\arabic*),leftmargin=2.2em,itemsep=2pt,topsep=2pt]
\item 剩余价值在生产领域中产生：资本家购买劳动力商品，劳动力的使用价值（劳动）能创造出大于自身价值的价值。\pingfen{3 分}
\item 工作日分为必要劳动时间和剩余劳动时间，剩余价值是工人在剩余劳动时间内创造的、被资本家无偿占有的价值。\pingfen{3 分}
\item 不变资本（c）只转移价值，不创造新价值；可变资本（v）购买劳动力，是剩余价值的唯一源泉。这一区分揭示了剥削的秘密。\pingfen{2 分}
\end{enumerate}

\thao{4.}\yaodian
\begin{enumerate}[label=(\arabic*),leftmargin=2.2em,itemsep=2pt,topsep=2pt]
\item 积极作用：促进生产要素在全球范围优化配置，推动技术扩散与产业转移，为发展中国家提供参与国际分工的机会。\pingfen{4 分}
\item 负面效应：加剧发展的不平衡，发达国家主导规则并获取大部分利益；增加金融风险传导；对发展中国家民族产业和生态环境形成压力。\pingfen{4 分}
\end{enumerate}

\datiA{论述题（每小题 15 分，共 30 分）}

\thao{1.}\yaodian
\begin{enumerate}[label=(\arabic*),leftmargin=2.2em,itemsep=2pt,topsep=2pt]
\item 社会基本矛盾是生产力与生产关系、经济基础与上层建筑的矛盾，它贯穿人类社会始终，规定社会形态的性质和更替。\pingfen{3 分}
\item 生产力是最活跃的因素，生产力的发展引起生产关系变革；生产关系反作用于生产力，适合时促进、不适合时阻碍。\pingfen{4 分}
\item 经济基础决定上层建筑，上层建筑反作用于经济基础；二者的矛盾运动推动社会形态由低级向高级演进。\pingfen{4 分}
\item 这一矛盾在阶级社会中表现为阶级矛盾与阶级斗争，在社会主义社会中则通过改革自我完善；因此改革是社会主义社会发展的直接动力。\pingfen{4 分}
\end{enumerate}

\thao{2.}\yaodian
\begin{enumerate}[label=(\arabic*),leftmargin=2.2em,itemsep=2pt,topsep=2pt]
\item \textbf{实质}：经济危机是生产相对过剩的危机，即相对于劳动群众有支付能力的需求而言的生产过剩，而非绝对过剩。\pingfen{4 分}
\item \textbf{根源}：生产社会化与生产资料资本主义私人占有之间的矛盾。个别企业生产的有组织性与整个社会生产的无政府状态的矛盾，生产无限扩大趋势与劳动群众有支付能力需求相对缩小的矛盾，都是这一基本矛盾的具体表现。\pingfen{5 分}
\item \textbf{无法根本解决}：资本主义的调整（国家干预、福利制度、金融化）只能暂时缓和矛盾、转移危机形式，不能改变私有制这一根源；只要基本矛盾存在，危机就会周期性爆发。\pingfen{4 分}
\item 结论：解决危机的根本出路在于以公有制取代私有制，这一判断是从基本矛盾推出的必然结论。\pingfen{2 分}
\end{enumerate}

\datiA{材料分析题（共 18 分）}

\thao{1.}
\textbf{（1）}\pingfen{9 分}
\begin{enumerate}[label=(\arabic*),leftmargin=2.2em,itemsep=2pt,topsep=2pt]
\item “两个必然”揭示的是历史发展的总趋势，回答的是资本主义向何处去的问题，属于规律层面的判断。\pingfen{3 分}
\item “两个决不会”揭示的是历史发展的条件性，回答的是变革何时发生的问题，强调生产力条件未成熟时旧制度不会灭亡。\pingfen{3 分}
\item 二者是方向与条件的统一：前者指明必然性，后者规定现实性；只讲“必然”易陷入空想，只讲“决不会”易陷入悲观等待。\pingfen{3 分}
\end{enumerate}

\textbf{（2）}\pingfen{9 分}
\begin{enumerate}[label=(\arabic*),leftmargin=2.2em,itemsep=2pt,topsep=2pt]
\item 社会主义代替资本主义是一个长期、曲折的历史过程，不能以某一时的挫折否定总趋势。\pingfen{3 分}
\item 必须把远大理想与现阶段任务统一起来：既坚持共产主义远大理想，又脚踏实地推进中国特色社会主义。\pingfen{3 分}
\item 当代资本主义仍有自我调节能力，其容纳生产力的空间尚未耗尽；这要求我们保持战略定力，把发展作为解决一切问题的基础。\pingfen{3 分}
\end{enumerate}
